\subsubsection{Physical vs Logical Gate Sets}

In quantum computing, it is important to distinguish between:
\begin{itemize}
    \item \definition{Physical Gate sets}, which correspond to the elementary \hl{operations directly implemented by the hardware};
    \item \definition{Logical Gate sets}, which are the \hl{fault-tolerant operations acting on encoded logical qubits} (e.g., the Clifford gates).
\end{itemize}
Although the two descriptions correspond to the same underlying quantum computation, they differ significantly in terms of implementation complexity and fault-tolerant requirements.

\highspace
\begin{flushleft}
	\textcolor{Green3}{\faIcon{microchip} \textbf{Physical (Hardware-Level) Gate Sets}}
\end{flushleft}
At the hardware level, quantum gates act directly on physical qubits. A typical universal physical gate set consists of:
\begin{itemize}
	\item \textbf{Single-Qubit Gates}: These include rotations around the X, Y, and Z axes (e.g., $R_x(\theta)$, $R_y(\theta)$, $R_z(\theta)$).

	In general, single-qubit rotations are \textbf{easier to implement} than multi-qubit gates because they \textbf{act locally on only one physical qubit}. In particular, many hardware platforms can implement $R_z (\theta)$ very efficiently, sometimes even virtually through software phase updates rather than physical operations.

	At the physical level, also the $T$ gate is simply a rotation around the Z-axis by $\pi/4$, therefore it is not more difficult to implement than other single-qubit gates.


	\item \textbf{Two-Qubit Gates}: The most common two-qubit gate is the CNOT gate, which is essential for creating entanglement between qubits.

	Two-qubit gates such as CNOT are \textbf{significantly more difficult} because they require \textbf{interactions between qubits}. In general, the \hl{larger the number of qubits involved in a gate, the more complex and error-prone its physical implementation becomes.}
\end{itemize}
These gates are \textbf{implemented through direct hardware control}, such as laser pulses in ion traps or microwave pulses in superconducting qubits.

\highspace
\begin{flushleft}
	\textcolor{Green3}{\faIcon{tools} \textbf{Logical (Fault-Tolerant) Gate Sets}}
\end{flushleft}
When quantum information is encoded into logical qubits using quantum error correction, the relevant gate set changes. The focus is no longer only on mathematical universality, but \textbf{also on fault-tolerant} implementability. A typical logical fault-tolerant gate set includes (\autopageref{sec:universal-set-of-quantum-gates}):
\begin{itemize}
	\item \textbf{Clifford gates}:
	\begin{equation*}
		\{H, S, \text{CNOT}\}
	\end{equation*}
	\item \textbf{Non-Clifford gates} (e.g., the $T$ gate):
	\begin{equation*}
		\{T\}
	\end{equation*}
\end{itemize}

\begin{flushleft}
	\textcolor{DarkOrange3}{\faIcon{balance-scale} \textbf{Key Distinction}}
\end{flushleft}
The crucial observation is that the \textbf{complexity of a gate depends strongly on whether it acts on}: physical qubits or encoded logical qubits. For example:
\begin{itemize}
	\item On physical qubits, $T = R_z(\pi/4)$ is a simple single-qubit rotation.
	\item On logical qubits, $T$ is a non-Clifford logical gate not usually transversally implementable, and therefore it requires magic state distillation and fault-tolerantly complex procedures to implement.
\end{itemize}
This distinction is one of the main practical challenges of large-scale quantum computing. At the mathematical level, the $T$ gate appears to be a simple phase rotation. However, once quantum information is encoded into logical qubits, implementing the same operation fault-tolerantly becomes extremely resource-intensive.

\highspace
As a consequence, the \hl{cost of fault-tolerant quantum algorithms is often dominated by:}
\begin{itemize}
	\item The number of logical $T$ gates required;
	\item Magic state preparation;
	\item Error-correction overheads associated with non-Clifford operations.
\end{itemize}
In summary, the same gate may be easy at the hardware level, but extremely difficult at the logical fault-tolerant level.